\documentclass[10pt]{article}
\usepackage[margin=0.8in]{geometry}
\usepackage{amsmath,amssymb,hyperref,booktabs}
\hypersetup{colorlinks=true,urlcolor=blue,linkcolor=blue}
\setlength{\parindent}{0pt}\setlength{\parskip}{4pt}
\newcommand{\F}{\mathbb F_5}
\begin{document}
\begin{center}
{\Large\bfseries Problem 00000002905: irreducibility does not force orbit divisibility}

lizaixi01 \quad AI-led research and formalization \quad 23 October 2026
\end{center}
\textbf{Result and scope.} We disprove the explicit divisibility clause for irreducible representations of finite linear groups. There is a faithful irreducible representation over $\F$, a generator whose matrix has order $2$, and a nonzero vector with a singleton orbit under that generator. Even its orbit under the whole group has odd cardinality $3$. Thus the prime $2$ divides the matrix order and divides neither actual orbit length.

The representation is irreducible for the whole group. Its restriction to the cyclic subgroup generated by the reflection is reducible; these are different properties. The source concerns generators of a linear group and irreducible representations of that group, and does not impose irreducibility on each cyclic restriction. We make no claim to refute that stronger hypothesis. Nor do we need to assign a meaning to the unspecified ``ratio-field'' tightness clause: the explicit divisibility condition already fails.

\section*{1. The actual finite group and matrices}
Let $G=D_3$ be the group of symmetries of a triangle. Write its elements as $r_i$ and $s_i$ for $i\in\mathbb Z/3\mathbb Z$, with
\[
r_i r_j=r_{i+j},\quad r_i s_j=s_{j-i},\quad
s_i r_j=s_{i+j},\quad s_i s_j=r_{j-i}.
\]
It has six elements and is generated by $r=r_1$ and $s=s_0$. The formal proof uses Mathlib's existing \texttt{DihedralGroup 3}, its cardinality theorem, and the actual subgroup closure of $\{r,s\}$. In particular, $5\nmid|G|$, so this is not a modular-characteristic example.

On $V=\F^2$ define
\[
R=\begin{pmatrix}0&-1\\1&-1\end{pmatrix},\qquad
S=\begin{pmatrix}0&1\\1&0\end{pmatrix},\qquad
\rho(r_i)=R^i,\quad \rho(s_i)=SR^i.
\]
The relations $R^3=I$, $S^2=I$, and $SRS=R^{-1}$ give a representation of $G$. Lean checks the full multiplication table as equalities of actual matrices, yielding a monoid homomorphism, then its invertible-matrix homomorphism into $\operatorname{GL}_2(\F)$. The six matrices are distinct, so the representation is faithful. The image is precisely the subgroup generated by $R$ and $S$, and its actual cardinality is $6$.

The determinants of $R$ and $S$ are $1$ and $-1$ respectively. Moreover $S\ne I$ and $S^2=I$, hence the matrix $S$ has actual order $2$. The formal theorem transfers the order of $s$ along the injective group homomorphism; it does not define the order to be $2$.

\section*{2. Irreducibility from a complete spanning argument}
Take any nonzero vector $u=(x,y)\in V$. Its image under $R$ is $Ru=(-y,x-y)$, and
\[
d(u)=\det[u\ Ru]=x^2-xy+y^2.
\]
This is nonzero for every nonzero $u$ over $\F$. One way to see this is that the discriminant of $X^2-X+1$ is $-3=2$, whereas the squares in $\F$ are $0,1,4$. The zero case $y=0$ is immediate. Lean directly checks the equivalent exhaustive statement
\[
\forall x,y\in\F,\qquad x^2-xy+y^2=0\ \Longrightarrow\ x=y=0.
\]
For every $w=(w_0,w_1)\in V$, the actual spanning identity is
\[
w=\frac{(x-y)w_0+yw_1}{d(u)}\,u+
  \frac{-yw_0+xw_1}{d(u)}\,Ru.
\]
Its denominator is proved nonzero before division. If a nonzero subrepresentation contains $u$, it contains $Ru$ and all these linear combinations, hence every $w$. Thus every subrepresentation is either $0$ or $V$, and $V\ne0$.

This establishes Mathlib's actual \texttt{Representation.IsIrreducible}: the lattice of actual invariant submodules has exactly its bottom and top. No abstract irreducibility hypothesis is assumed, and the proof covers all invariant subspaces, not a selected list.

\section*{3. Both genuine orbit cardinalities}
Let $v=(1,1)$, which is nonzero. The reflection fixes it: $Sv=v$. Every power of $S$ therefore fixes $v$, and its cyclic-subgroup orbit is exactly
\[
\langle s\rangle v=\{(1,1)\},\qquad |\langle s\rangle v|=1.
\]
The six actual group elements give the following images, with entries represented modulo $5$:
\begin{center}
\begin{tabular}{c|cccccc}
\toprule
$g$ & $r_0$ & $r_1$ & $r_2$ & $s_0$ & $s_1$ & $s_2$\\
\midrule
$\rho(g)v$ & $(1,1)$ & $(4,0)$ & $(0,4)$ & $(1,1)$ & $(0,4)$ & $(4,0)$\\
\bottomrule
\end{tabular}
\end{center}
Consequently
\[
Gv=\{(1,1),(4,0),(0,4)\},\qquad |Gv|=3.
\]
Lean uses the existing group-action orbit sets, proves these set equalities, proves finiteness, and evaluates their actual cardinalities. Neither $1$ nor $3$ is supplied as the definition of an orbit length. The cyclic orbit proof uses the actual subgroup of integer powers and its inclusion in the vector stabilizer.

Since $2$ is prime and divides $\operatorname{ord}(S)=2$, while
\[
2\nmid|\langle s\rangle v|\quad\hbox{and}\quad 2\nmid|Gv|,
\]
the stated divisibility fails whether the orbit in that clause means the generator orbit or the whole-group orbit. The two objects are separately defined and computed, rather than interchanged.

\section*{4. Verification and limits}
The complete original bilingual source is in \texttt{SOURCE.md}. The pinned Lean project contains the full proof with an audit of all 53 named declarations: 35 theorems, 13 definitions, 3 abbreviations, and 2 instances. This includes the representation, group action, matrix homomorphisms, and irreducibility construction, as well as all orbit statements. Finite table checks use ordinary kernel-evaluated \texttt{decide}; there is no \texttt{native\_decide}, \texttt{sorry}, or extra axiom.

The submission includes source-bound build, axiom, and full fresh-kernel receipts. Worker and separate supervising AI review the complete mathematical scope and rendered PDF pages. No organizer acceptance is claimed. No earlier same-ID submission was found in the eligibility checks.

The result concerns whole-group irreducibility, and does not dispute statements assuming that the chosen generator itself acts irreducibly on $V$. It also does not claim a classification of ratio fields or eigenvalue configurations. The counterexample is over a genuine finite field, is faithful, has nonzero vector and nonmodular group order, and refutes only the necessary divisibility clause. This is selective production, not a new controlled efficiency measurement.
\end{document}
